\documentclass[12pt,a4paper]{article}
\usepackage{polyglossia}
\setmainlanguage{hebrew}
\setotherlanguage{english}
\usepackage{fontspec}
\newfontfamily\hebrewfont[Script=Hebrew]{David CLM}
\newfontfamily\hebrewfontsf[Script=Hebrew]{Miriam CLM}
\setmainfont[Script=Hebrew]{David CLM}
\usepackage{amsmath}
\usepackage{amssymb}
\usepackage{xcolor}
\usepackage[margin=1in]{geometry}
\usepackage{enumitem}
\usepackage{parskip}
\usepackage{titlesec}
\usepackage{microtype}

\titleformat{\section}{\large\bfseries}{}{0em}{}[\titlerule]
\titleformat{\subsection}{\normalsize\bfseries\itshape}{}{0em}{}

\title{לוגיקה למדעי המחשב}
\author{}
\date{}

\begin{document}
\maketitle

%---------------------------------------------------------------------
\section*{הרצאה 1:}
%---------------------------------------------------------------------

\begin{quote}
\textenglish{``If I could begin to be half of what you think of me, I could do about anything\ldots\ I could even learn how to love''}
\end{quote}

\ldots\ (תכונות ל $X_{B,F}$)

%---------------------------------------------------------------------
\section*{הרצאה 2:}
%---------------------------------------------------------------------

\textenglish{``I gave my trust, I shed my blood. I held on tight just to end up stung\ldots\ you were my friend- scorpion!''}

בהנתן עולם $W$ ובסיס $B$ וקבוצת פונקציות מעל $W$, הקבוצה $X_{B,F}$ קיימת ויחידה.
\textcolor{red}{כלומר כאשר מתארים קבוצה אינדוקטיבית מספיק להצביע על השלשה הזו.}

איך מראים שקבוצה היא $X_{B,F}$? מראים לפי ההגדרה\ldots\ או משתמשים בהגדרות השקולות הבאות:

$X_{B,F}=\bigcap\{T\subset W \mid B\subset T \text{ and } T \text{ is closed under } W\}$

נגדיר: $D_{0}=B,\ D_{1}=D_{0}\cup F(D_{0}),\ldots,\ D_{i}=D_{i-1}\cup F(D_{i-1}),\ldots$, עתה: $X_{B,F}=\bigcup_{i\in\mathbb{N}}D_{i}$

\subsection*{תחשיב הפסוקים:}

נגדיר קבוצות:

$Var=\{p_{i} \mid i\in\mathbb{N}\}$

$Symb=\{\wedge,\vee,\rightarrow,\lnot,(,)\}$

$W=(Var\cup Symb)^{*}$

$B=Var$

$F=\{F_{\lnot},\ldots\}$

$WFF=X_{B,F}$ תיקרא תחשיב הפסוקים.

בפרט:

אם רוצים להראות שתחשיב הפסוקים אמ''מ משהו (תכונה) יש להראות לפי ההגדרה או השקילויות.

אם רוצים להראות שתחשיב הפסוקים מקיימת תכונה מסויימת צריך להראות את זה על ידי אינדוקציית מבנה.

אם רוצים להראות שאיבר שייך לתחשיב הפסוקים צריך להראות סדרת יצירה לאיבר!

ועוד\ldots

\subsection*{משפט הקריאה היחידה}

\ldots

\subsection*{סמנטיקה של תחשיב הפסוקים:}

נגדיר השמה $z:Var\rightarrow\{0,1\}$.
$Ass=\{z \mid z:Var\rightarrow\{0,1\}\}$

%---------------------------------------------------------------------
\section*{הרצאה 3:}
%---------------------------------------------------------------------

\begin{quotation}
\textenglish{``There's no time, we're losing the light''}
\end{quotation}

נגדיר השמה מורחבת $\overline{z}$ מעל $WFF$ בצורה אינדוקטיבית:

לכל $v\in Var$: $\exists i\in\mathbb{N},\ \overline{z}(v)=z(v)=z(p_{i})$

סגור: \ldots

משפט התלות הסופית: בהנתן פסוק $\alpha$ ושתי השמות $z_{1},z_{2}$ שמקיימות $z_{1}(p_{i})=z_{2}(p_{i})$ לכל אטום המופיע בפסוק, אזי $z_{1}\vDash\alpha\iff z_{2}\vDash\alpha$.

בצורה הזו, פסוק מתאר טבלת אמת.

\emph{האם לכל טבלת אמת קיים פסוק שמתאר אותה?}

\textcolor{blue}{לא! אבל תחת תנאים מסויימים, כן.}

\textbf{הגדרה:} מערכת קשרים שעבורה מתקיים שלכל טבלת אמת יש פסוק שמתאר אותה נקראת מערכת קשרים שלמה.

\textcolor{red}{יש לחזור על הדוגמאות מההרצאה + הוכחות שמערכת קשרים היא שלמה או לא.}

%---------------------------------------------------------------------
\section*{הרצאה 4:}
%---------------------------------------------------------------------

\begin{quote}
``אם ביחד, אז רק לחוד''
\end{quote}

פסוק $\alpha$ נקרא טאוטולוגיה או אמת לוגית אם לכל השמה $z$ מתקיים $z\vDash\alpha$.

\textbf{איך נראה שפסוק הוא טאואולוגיה? טבלת אמת (משפט התלות הסופית) או הנחה בשלילה.}

\underline{טענה:} $\alpha$ סתירה אמ''מ $\lnot\alpha$ טאוטולוגיה. $\alpha$ טאוטולוגיה אמ''מ $\lnot\alpha$ סתירה.

דוגמאות חשובות לטאוטולוגיות:

$\alpha\rightarrow(\beta\rightarrow\alpha)$

$\alpha\rightarrow(\beta\rightarrow\gamma)\rightarrow((\alpha\rightarrow\beta)\rightarrow(\alpha\rightarrow\gamma))$

$(\lnot\beta\rightarrow\lnot\alpha)\rightarrow(\alpha\rightarrow\beta)$

\subsubsection*{נביעה לוגית:}

נאמר שפסוק $\beta$ נובע לוגית מפסוק $\alpha$ אם לכל השמה $z$ מתקיים: אם $z\vDash\alpha$ אז $z\vDash\beta$. ונסמן $\alpha\vDash\beta$.

\underline{טענה:} $\alpha\vDash\beta$ אמ''מ $\vDash\alpha\rightarrow\beta$.

\subsubsection*{שקילות לוגית}

נאמר ששני פסוקים $\alpha,\beta$ שקולים לוגית אם $\alpha\vDash\beta$ וגם $\beta\vDash\alpha$.

אינטואיציה: פסוקים שקולים הם פסוקים שמספרים ``אותו הסיפור''. בפרט, כל פסוק שקול לעצמו.

\subsubsection*{הרחבת המושגים לקבוצות של פסוקים:}

קבוצת פסוקים $\Sigma$ תיקרא ספיקה אם: $\exists z,\forall\alpha\in\Sigma,\ z\vDash\alpha$. קבוצת פסוקים יכולה להיות אינסופית ויכולה להיות ריקה. בפרט עבור קבוצת פסוקים ריקה היא ספיקה, וכל ההשמות ($Ass$) מספקות אותה.

\textbf{אינטואיציה: $\land$ על כל הפסוקים. ונותן לנו דרך לבטא את זה במקרה האינסופי. בפרט עבור קבוצה סופית ניתן לבטא אותה על ידי $\underset{\alpha\in X}{\bigwedge}\alpha$ (ורק עבור סופית!). אם הקבוצה ריקה אז הביטוי לא מוגדר.}

יש גם הגדרות לשקילות קבוצות פסוקים ולגרירה לוגית לפסוק מקבוצה.

\subsubsection*{טענות נוספות:}

\begin{enumerate}
  \item $X\cup\{\alpha\}\vDash\beta$ אמ''מ $X\vDash\alpha\rightarrow\beta$.
  \item אם $X\cup\{\alpha\}\vDash\beta$ וגם $X\vDash\alpha$ אז $X\vDash\beta$.
  \item אם $X\cup\{\alpha\}\vDash\beta$ וגם $X\cup\{\lnot\alpha\}\vDash\beta$ אזי $X\vDash\beta$.
  \item אם $X\vDash\beta$ וגם $X\vDash\lnot\beta$ אז $X$ לא ספיקה.
  \item אם $X\cup\{\lnot\alpha\}\vDash\beta$ וגם $X\cup\{\lnot\alpha\}\vDash\lnot\beta$ אז $X\vDash\alpha$.
\end{enumerate}

ויש עוד\ldots\ (3)

\underline{הערה (1):} קבוצת ההשמות המספקות קבוצה $\Sigma$ מסומנת על ידי $M(\Sigma)$.

\underline{הערה (2):} הרצאה מעניינת :(

%---------------------------------------------------------------------
\section*{הרצאה 5:}
%---------------------------------------------------------------------

\begin{quote}
\textenglish{``And is all that I might owe you carved on ivory?''}
\end{quote}

\subsection*{מערכת הוכחה שלמה}

\underline{הגדרה:} מעל העולם $WFF_{\{\lnot,\rightarrow\}}$ וקבוצת פעולות $F=\{MP\}$ וקבוצת בסיס $B$ נגדיר קבוצה אינדוקטיבית בצורה הבאה:

בסיס $B$: אקסיומות

יהיו $\alpha,\beta,\gamma\in WFF_{\{\lnot,\rightarrow\}}$:

$\alpha\rightarrow(\beta\rightarrow\alpha)\in B$

$\alpha\rightarrow(\beta\rightarrow\gamma)\rightarrow((\alpha\rightarrow\beta)\rightarrow(\alpha\rightarrow\gamma))\in B$

$(\lnot\beta\rightarrow\lnot\alpha)\rightarrow(\alpha\rightarrow\beta)\in B$

פעולות (כללי היסק): יהיו $\alpha,\beta,\gamma\in X_{B,F}$. אזי:

$MP(\alpha,\alpha\rightarrow\beta)=\beta$

אם $\gamma\neq\alpha\rightarrow\beta$ אזי $MP(\alpha,\gamma)=\alpha$

סימון היסק: $MP:\dfrac{\alpha,\;\alpha\rightarrow\beta}{\beta}$

קבוצה זו נקראת קבוצת הפסוקים היכיחים. כאשר לכל $\alpha\in X_{B,F}$ נסמן $\vdash\alpha$.

\subsubsection*{הוכחה מתוך הנחות}

בהנתן קבוצת פסוקים $\Sigma$, הקבוצה $X_{B,F}$ תכלול גם את כל הפסוקים שב-$\Sigma$ בבסיס שלה.

סימון $Ded(X):=X_{B,F}$.

\textbf{תכונות של הוכחה מתוך הנחות:}

\begin{enumerate}
  \item אם $\alpha\in X$ אז $X\vdash\alpha$
  \item אם $X\subseteq Y$ וגם $X\vdash\alpha$ אז $Y\vdash\alpha$ (מונוטוניות)
  \item אם לכל $\beta\in X$ מתקיים $Y\vdash\beta$ אזי לכל $X\vdash\alpha$ מתקיים $Y\vdash\alpha$.
\end{enumerate}

\underline{הוכחה ל-3:}

נבנה סדרת יצירה ל-$\alpha$ מתוך $Y$.

יהי $X\vdash\alpha$. אזי קיימת סדרת הוכחה ל

\noindent 1.

\noindent 2.

$\vdots$

\noindent $k$.

משפט הדדוקציה:

לכל $\alpha,\beta\in WFF_{\{\lnot,\rightarrow\}}$ ולכל $X\subseteq WFF_{\{\lnot,\rightarrow\}}$ מתקיים:

$X\vdash\alpha\rightarrow\beta$ אמ''מ $X\cup\{\alpha\}\vdash\beta$.

אינטואיציה: לפעמים יותר נוח לצרף את $\alpha$ לקבוצת ההנחות שלנו על מנת להוכיח משהו.

\textbf{איך מוכיחים שפסוק יכיח מתוך קבוצה?}

לפי משפט הדדוקציה ואז תכונות, או על פי יצירת סדרת הוכחה.

\subsubsection*{משפט הנאותות:}

לכל פסוק $\alpha$ ולכל קבוצת פסוקים $X$ מתקיים:

$X\vDash\alpha\Leftarrow X\vdash\alpha$.

\textbf{הוכחה: אינדוקציית מבנה}

\subsubsection*{קבוצת פסוקים עקבית}

יש שתי הגדרות שקולות:

\begin{enumerate}
  \item קבוצת פסוקים $X$ נקראת עקבית אם קיים פסוק $\alpha$ כך ש-$X\nvdash\alpha$.
  \item קבוצת פסוקים $X$ נקראת עקבית אם לכל פסוק $\alpha$ מתקיים $X\vdash\alpha$ או $X\nvdash\lnot\alpha$.
\end{enumerate}

\textcolor{red}{להסתכל על ההוכחה!}

\textbf{משפט}: אם $X$ ספיקה אז היא עקבית.

הערות:

\begin{itemize}
  \item \textcolor{blue}{כדאי לחשוב האם כללי ההיסק השונים מקיימים את משפט הנאותות $(Ded(X)\subseteq Con(X))$. למשל: $\dfrac{\alpha\wedge\beta}{\alpha,\beta},\ \dfrac{\alpha}{\lnot\lnot\alpha},\ \dfrac{\alpha\rightarrow\beta,\lnot\beta}{\lnot\alpha},\ \dfrac{\lnot(\alpha\land\beta)}{\lnot\alpha,\lnot\beta}$.}
\end{itemize}

\subsubsection*{תכונות מערכת ההוכחה (תרגול):}

\begin{itemize}
  \item הנחת המבוקש: $\Sigma\vdash\alpha\Leftarrow\alpha\in\Sigma$
  \item סופיות ההוכחה: $\exists\Sigma:|\Sigma|<\infty\wedge\Sigma\subseteq\Gamma,\ \Sigma\vdash\alpha\Leftarrow\Gamma\vdash\alpha$
  \item מונוטוניות (ראינו)
\end{itemize}

%---------------------------------------------------------------------
\section*{הרצאה 6:}
%---------------------------------------------------------------------

\begin{quote}
\textenglish{``It works! Finally, it works! I don't know why, but it does. And that's all that matters''}
\end{quote}

\textbf{משפט:}

תהי $X$ קבוצת פסוקים ויהי $\alpha$ פסוק.

אם $X\cup\{\lnot\alpha\}$ לא עקבית אז $X\vdash\alpha$.

\textcolor{red}{להסתכל על ההוכחה של זה!}

\textbf{משפט:}

תהי $X$ קבוצת פסוקים. אם $X$ עקבית אז $X$ ספיקה.

\textbf{משפט השלמות:}

תהי $X$ קבוצת פסוקים ויהי $\alpha$ פסוק. $X\vdash\alpha\Leftarrow X\vDash\alpha$.

\subsubsection*{עקביות מקסימלית}

קבוצת פסוקים $Y$ נקראת עקבית מקסימלית אם היא עקבית וגם לכל פסוק $\alpha$ מתקיים $X\vdash\alpha$ או $X\vdash\lnot\alpha$.

משפטים:

\begin{itemize}
  \item אם $X$ עקבית אז קיימת $Y$ עקבית מקסימלית כך ש-$X\subseteq Y$: \textcolor{red}{הוכחה חשובה}
  \item $X$ עקבית מקסימלית אמ''מ $|M(X)|=1$
\end{itemize}

\subsubsection*{משפט הקומפקטיות}

$Y$ עקבית מקסימלית אמ''מ כל תת-קבוצה \textbf{סופית} שלה היא עקבית.

\subsubsection*{טיפים להוכחות מהצורה $Ded_{N}(\emptyset)\subseteq Con(\emptyset)$:}

\begin{itemize}
  \item וודא/י שמתקיים עבור כל כללי ההיסק שהפלט הוא אכן טאוטולוגיה בהינתן שהקלט טאוטולוגיה, וגם עבור האקסיומות שהן אכן טאוטולוגיות.
  \begin{itemize}
    \item אם זהו אכן המצב, תמשיך/י להוכחה באינדוקציית מבנה ``רגילה'' וסיימת/ה.
  \end{itemize}
  \item אחרת, וודא/י שכללי ההיסק שלא מקיימים את זאת אכן ניתנים לשימוש, על ידי למשל בדיקת האם ניתן בכלל להגיע לביטוי מהצורה של הארגומנטים (קשר מרכזי וכו').
  \begin{itemize}
    \item אם זהו אכן המצב, הגדר/י תכונה $T$ מתאימה המקיימת $T\subseteq Con(\emptyset)$ והראה/י שמתקיים $Ded_{N}(\emptyset)\subseteq T$.
  \end{itemize}
  \item אחרת, כנראה יש להפריך\ldots\ ואז בהצלחה לכולנו :(
\end{itemize}

%---------------------------------------------------------------------
\section*{הרצאה 7:}
%---------------------------------------------------------------------

\begin{quote}
\textenglish{``If I lead you through the fury, will you call to me?''}
\end{quote}

\subsection*{גדירות}

נאמר שקבוצת ההשמות $K\subseteq Ass$ היא גדירה אם קיימת קבוצת פסוקים $\Sigma$ שמגדירה אותה. כלומר קיימת $\Sigma$ כך ש-$M(\Sigma)=K$.

\subsection*{גדירות באופן סופי}

נאמר שקבוצת ההשמות $K$ היא גדירה באופן סופי אם קיימת $\Sigma$ סופית שמגדירה אותה.

\subsubsection*{משפט: התנאים הבאים שקולים}

\begin{itemize}
  \item $K$ גדירה באופן סופי
  \item $K$ גדירה על ידי פסוק יחיד
  \item $K,\ Ass\setminus K$ גדירות
\end{itemize}

%---------------------------------------------------------------------
\section*{הרצאה 8:}
%---------------------------------------------------------------------

\begin{quote}
\textenglish{``Jesse, I don't know what I have done''}
\end{quote}

\subsection*{תחשיב היחסים}

א''ב של השפה:

סימנים שאינם תלויים במילון:

\begin{itemize}
  \item $Symb$
  \item כמתים
  \item סימן יחס שוויון $\thickapprox$
  \item $Var$
\end{itemize}

סימנים שכן תלויים במילון:

\begin{itemize}
  \item סימני קבוע: $c_{i}$ עבור $i\in\mathbb{N}$
  \item סימני פונקציה: $F_{i,k}$ עבור $i\in\mathbb{N}$, $k\geq1$ סימן פונקציה $k$-מקומית
  \item סימני יחסים: $R_{i,k}$ עבור $i\in\mathbb{N}$, $k\geq1$ סימן יחס $k$-מקומי
\end{itemize}

\underline{שורה תחתונה:} על מנת להגדיר מילון מספיק לציין סימני קבועים, פונקציות ויחסים.

\subsection*{שמות עצם:}

קבוצה אינדוקטיבית המוגדרת מעל קבוצת המחרוזות הסופיות מעל הא''ב מעל $\tau$.

\textbf{בסיס:}

משתנים $Var$ וסימני קבוע שמופיעים ב-$\tau$.

\textbf{פעולות:}

לכל סימן פונקציה $k$-מקומית $F$ שמופיע במילון מוגדרת פעולה $W^{k}\rightarrow W$ שבהנתן $t_{1},t_{2},\ldots,t_{k}\in W$ מחזירה את $F(t_{1},t_{2},\ldots,t_{k})$.

\subsubsection*{משפט הקריאה היחידה לשמות עצם}

לכל שם עצם ``יש פירוק יחיד''.

\underline{הערה מהמחברת 1}: בהרצאה הרלוונטית לא הוצג בצורה פורמלית\ldots\ \textenglish{so don't come for me}

\underline{הערה מהמחברת 2}: זה יהיה המצב גם בהמשך כשאציג אותו עבור קבוצת הנוסחאות :>

\subsection*{נוסחאות:}

קבוצה אינדוקטיבית המוגדרת מעל קבוצת המחרוזות הסופיות מעל הא''ב מעל $\tau$.

\textbf{בסיס:}

לכל שני שמות עצם מעל $\tau$ שנסמן אותם על ידי $t_{1},t_{2}$ מתקיים ש-$t_{1}\approx t_{2}$ היא נוסחה.

לכל סימן יחס $k$-מקומי $R$ ב-$\tau$ ולכל $t_{1},t_{2},\ldots,t_{k}$ שמות עצם מעל $\tau$:
$R(t_{1},t_{2},\ldots,t_{k})$ היא נוסחה.

\underline{אלו ייקראו נוסחאות אטומיות.}

\textbf{פעולות:}

\begin{itemize}
  \item הפעלת קשרים
  \item הפעלת כמתים
  \begin{itemize}
    \item \underline{הערה}: הקפידו על סוגריים!
  \end{itemize}
\end{itemize}

\subsubsection*{משפט הקריאה היחידה לנוסחאות:}

לכל נוסחה ``יש פירוק יחיד''.

\subsection*{סמנטיקה לתחשיב היחסים}

בהינתן מילון $\tau=\langle c_{0},c_{1},\ldots,F_{0},F_{1},\ldots,R_{0},R_{1}\rangle$, מבנה עבור המילון ייראה כך:
$M=\langle D^{M},c_{0}^{M},c_{1}^{M},\ldots,F_{0}^{M},F_{1}^{M},\ldots,R_{0}^{M},R_{1}^{M}\rangle$.

כאשר:

\begin{itemize}
  \item $D^{M}\neq\emptyset$
  \item לכל סימן קבוע $c_{i}$ ב-$\tau$, $c_{i}^{M}\in D^{M}$
  \item לכל פונקציה \ldots
\end{itemize}

\subsection*{הרצאה 9:}

\begin{quote}
\textenglish{``when I ask a good question, people assume I'm a guy. When I ask a stupid question, people assume I'm a girl.''}
\end{quote}

\subsection*{השמות}

בהינתן מילון $\tau$ ומבנה $M$ מעל המילון, השמה היא פונקציה $z:\{v_{i} \mid i\in\mathbb{N}\}\rightarrow D^{M}$.

\textbf{השמה מורחבת:} $\overline{z}:term(\tau)\rightarrow D^{M}$. מוגדרת בצורה הבאה:

בסיס: $\overline{z}(t)=z(t)=z(v_{i})$, עבור קבועים פשוט נותנת להם את הקבוע.

צעד אינדוקציה: $t=F(t_{1},\ldots,t_{k})$, $\overline{z}(t)=F^{M}(\overline{z}(t_{1}),\ldots,\overline{z}(t_{k}))$

\subsection*{ספיקות}

``מתי מבנה $M$ מעל מילון $\tau$ והשמה $z$ מספקים נוסחא?''

\subsection*{משתנים קשורים וחופשיים}

משתנה ייקרא משתנה חופשי אם קיים מופע חופשי שלו.

\subsection*{משפט אנאלוגי למשפט התלות הסופית}

יהי $\tau$ מילון ותהי נוסחה $\alpha$ מעל $\tau$.

יהי $M$ מבנה עבור $\tau$ ותהיינה $z_{1},z_{2}$ השמות עבור $M$.

אם מתקיים $z_{1}(v_{i})=z_{2}(v_{i})$ לכל משתנה $v_{i}$ חופשי ב-$\alpha$, אז $z_{2}$ מספקת את $\alpha$ אמ''מ $z_{1}$ מספקת אותה.

נוסחה ללא משתנים חופשיים נקראת פסוק.

פסוק לא תלוי בערך ש-$z$ נותנת למשתנים.

%---------------------------------------------------------------------
\part*{הרצאה 10:}
%---------------------------------------------------------------------

\begin{quote}
\textenglish{``We haven't had our pizza party yet :(''}
\end{quote}

\subsection*{ספיקות של נוסחה}

בהינתן מילון $\tau$.

נוסחה $\alpha$ תיקרא ספיקה אם קיים מבנה $M$ והשמה $z$ כך ש-$M\vDash_{z}\alpha$.

אם קיים מבנה $M$ המקיים שלכל $z$, $M\vDash_{z}\alpha$, נסמן $M\vDash\alpha$.

אם לכל מבנה $M$ מתקיים שלכל $z$, $M\vDash\alpha$, נסמן $\vDash\alpha$.

נוסחה $\alpha$ נקראת סתירה אם לכל מבנה $M$ ולכל $z$\ldots

\subsection*{גרירה לוגית}

בהינתן מילון $\tau$, נאמר שקבוצת נוסחאות $\Sigma$ גוררת לוגית $\alpha$ אם לכל מבנה $M$ והשמה $z$ מתקיים שאם $M\vDash_{z}\Sigma$ אז $M\vDash_{z}\alpha$.

במקרה זה נסמן $\Sigma\vDash\alpha$.

אם $\Sigma=\{\beta\}$ ומתקיים $\Sigma\vDash\alpha$ אז נסמן $\beta\vDash\alpha$.

אם $\alpha\vDash\beta$ וגם $\beta\vDash\alpha$ אז נאמר שהם שקולים לוגית ונסמן $\alpha\equiv\beta$.

\textcolor{red}{לעשות דוגמה אחת מכל סוג.}

\subsection*{גדירות}

בהינתן $\tau$ עבור קבוצת פסוקים $\Sigma$ נגדיר $M(\Sigma)=\{M \mid M\vDash\Sigma\}$.

נאמר שקבוצת מבנים $K$ היא גדירה אם קיימת קבוצת נוסחאות $\Sigma$ כך ש-$M(\Sigma)=K$.

\subsection*{משפט הקומפקטיות}

קבוצת נוסחאות סגורות נקראת ספיקה אמ''מ לכל תת-קבוצה סופית מתקיים שהיא ספיקה.

על מנת להפריך גדירות נשתמש במשפט הזה.

%---------------------------------------------------------------------
\part*{הרצאה 11:}
%---------------------------------------------------------------------

\begin{quote}
\textenglish{``Darker than your brown?''}
\end{quote}

%---------------------------------------------------------------------
\part*{הרצאה 12:}
%---------------------------------------------------------------------

\begin{quote}
\textenglish{``Are alpha women a thing?''}
\end{quote}

\end{document}
